\chapter{Conclusions and Future Work}
\label{chap:conclusions}

This report investigated the modernization of Outsight Cloud through the migration of a production platform from a legacy infrastructure toward a cloud-native architecture on AWS. The original system relied on self-managed infrastructure and Docker Compose while depending on several external services, most notably Airtable as the primary operational datastore. Although this enabled rapid early development, it progressively introduced limitations related to scalability, maintainability, operational governance, and dependency management. Rather than approaching migration as a simple infrastructure relocation, this work treated cloud adoption as a broader architectural transformation involving infrastructure modernization, data-layer evolution, deployment automation, and operational redesign, supported by concrete engineering artifacts (Terraform definitions, data migration tooling, validation scripts, and environment-specific deployment workflows).

\section{Key Outcomes}
The evaluation in Chapter~\ref{chap:evaluation} provides quantitative evidence across three dimensions. From a cost perspective (Section~\ref{sec:costanalysis}), the migration delivers a structural saving of €235/month driven primarily by replacing Airtable with Amazon RDS, eliminating a per-seat SaaS pricing model ill-suited to a backend role, reinforced by reduced object-storage costs. From an operational standpoint (Section~\ref{sec:operational}), delegating patching, scaling, deployment, and database administration to managed services reduces the team's undifferentiated operational burden by roughly 90\,\%, both in hours and cost. From a performance perspective, the HTTP benchmark shows consistent improvements across all timing components, most pronounced in the DNS and TCP phases, with a tighter latency distribution that brings worst-case user experience closer to the median.

Beyond the measured outcomes, two structural contributions complete the picture. The adoption of Terraform established a reproducible, version-controlled, and auditable provisioning baseline, re-establishing ownership over infrastructure and data assets. The incremental migration design, supported by dual-backend coexistence, data replication tooling, validation scripts, and environment separation, minimized migration risk and service disruption while establishing a foundation of independently scalable managed services capable of supporting long-term platform evolution.

\section{Limitations and Open Questions}
Two aspects of the measurement scope should be acknowledged. The latency benchmark targets the publicly accessible homepage endpoint, which provides a reproducible, infrastructure-representative baseline; authenticated API endpoints and database-intensive workflows were not included, as their evaluation requires a stable authentication token and a consistent asset fingerprint across both environments. The operational effort figures are based on informed engineering estimates grounded in direct team experience rather than systematically recorded timesheet data.

\section{Future Work}
Several directions remain open. The benchmark coverage could be extended to authenticated API endpoints and database query performance, complementing the infrastructure-level measurements with application-layer evidence. Load testing under concurrent demand would confirm whether ECS Fargate auto-scaling and RDS connection pooling perform as expected under peak conditions. The identity and access management layer introduced through Amazon Cognito could be further standardized, and observability could be extended through richer CloudWatch dashboards and post-cutover performance baselines. Finally, the migration principles documented here (the decision framework, IaC-first provisioning, dual-backend validation, and incremental rollout planning) could be evaluated in other industrial contexts to assess their generalizability.

Overall, this report demonstrates that legacy-to-cloud-native migration is most effective when treated as a structured engineering process supported by explicit architectural decisions, reproducible infrastructure management, continuous validation, and controlled incremental evolution. The Outsight Cloud case study shows how modernization can be grounded in concrete engineering practices, and the quantitative evaluation confirms measurable improvements across cost, latency, and operational overhead, validating the architectural decisions made throughout the process.
